\documentclass[12pt]{article}
\usepackage[utf8]{inputenc}
\usepackage[T1]{fontenc}
\usepackage{hyperref}
\title{Tunggare Language Resources for K-12 Curriculum Development}
\author{Compiled by Claude Sonnet 5 (claude-sonnet-5) for Andr\'as Kornai}
\date{}
\begin{document}
\maketitle

\section{Introduction}
Tunggare (ISO 639-3 trt; also spelled Tarunggare or Turunggare) is a small East Geelvink Bay language of Papua, Indonesia, near the Nabire coast. Search found only encyclopedic and people-group entries, with no educational or documentation materials.

\section{Primary Educational Resources}
None found.

\section{Linguistic and Reference Materials}
\begin{itemize}
\item Wikipedia entry on Tunggare: \url{https://en.wikipedia.org/wiki/Tunggare_language}
\item Wikipedia entry on Tarunggare (near-duplicate): \url{https://en.wikipedia.org/wiki/Tarunggare_language}
\item Wikipedia on the East Geelvink Bay family: \url{https://en.wikipedia.org/wiki/East_Geelvink_Bay_languages}
\end{itemize}

\section{Religious and Community Resources}
Joshua Project people-group profile (speaker figures differ from Wikipedia): \url{https://joshuaproject.net/index.php/people_groups/15276}

\section{Digital and Audio Resources}
None found.

\section{Curriculum Development Context}
No grammar, dictionary, orthography or school materials were found. Any curriculum would need fieldwork with the community first.
\end{document}